% GENERATED FILE. Edit canonical lessons, then run npm run generate:books.
% canonical-source-hash: fnv1a64:e5a449bdecf4921d
% canonical-lessons: GE-C505-trainieren, GE-C505-ausdauer, GE-C505-kondition, GE-C505-dehnen, GE-C505-muskelkater, GE-R505-review-498-505

\chapter{Trainieren}
\label{ch:ge-trainieren}
\hlchaptermodality{505}

\begin{chapteropening}
\textbf{By the end of this chapter:} \emph{I can talk about training and fitness.}

Complete the last lesson of chapter 505: I can talk about training and fitness.
\end{chapteropening}

\section[trainieren]{trainieren — to train, to practise}

\label{lesson:GE-C505-trainieren}



\begin{quote}
Before the new one: say the German for seed, then the German for harvest.
\end{quote}

\subsection*{You'll want to know: trainieren}
\textbf{trainieren} — \textquotedblleft{}to train, to practise\textquotedblright{}.

\begin{quote}
Die Mannschaft trainiert dreimal pro Woche.
\end{quote}

\begin{quote}
\emph{The team trains three times a week.}
\end{quote}

Regular, Perfekt without \emph{ge-}: \emph{hat trainiert}.

\begin{cousinweb}
\emph{trainieren} was borrowed from English \emph{to train} (EtymDB lists the German verb), which came from French \emph{traîner} \textquotedblleft{}to drag\textquotedblright{}, from Latin \emph{trahere} \textquotedblleft{}to pull\textquotedblright{}.
\end{cousinweb}

\begin{grammarlens}[title={What you've built}]
One more word for telling, explaining and describing at B1.
\end{grammarlens}

\subsection*{Your turn}
\begin{itemize}
\raggedright
  \item \emph{Say it:} \emph{trainieren}
  \item \emph{Say it:} \emph{trainieren}, once more
  \item \emph{Say it:} a sentence of your own with \emph{trainieren}
  \item \emph{Recall:} say the German for seed, then the German for harvest, then say \emph{trainieren} again
\end{itemize}

\begin{tcolorbox}[breakable,colback=teal!4,colframe=teal!35!black,title={Before you move on}]
What does trainieren mean? (\textquotedblleft{}To train, to practise\textquotedblright{}.) Say it once more.
\end{tcolorbox}
\section[die Ausdauer]{die Ausdauer — stamina, endurance}

\label{lesson:GE-C505-ausdauer}



\begin{quote}
Before the new one: say the German for harvest, then the German for to train.
\end{quote}

\subsection*{You'll want to know: die Ausdauer}
\textbf{die Ausdauer} — \textquotedblleft{}stamina, endurance\textquotedblright{}.

\begin{quote}
Für einen Marathon braucht man viel Ausdauer.
\end{quote}

\begin{quote}
\emph{For a marathon you need a lot of stamina.}
\end{quote}

No plural.

\begin{cousinweb}
\emph{Ausdauer} is built on \emph{aus} \textquotedblleft{}out\textquotedblright{} and \emph{dauern} \textquotedblleft{}to last\textquotedblright{}, from Middle High German \emph{dūren}, which came from Latin \emph{dūrāre} \textquotedblleft{}to last\textquotedblright{} (EtymDB).
\end{cousinweb}

\begin{grammarlens}[title={What you've built}]
One more word for telling, explaining and describing at B1.
\end{grammarlens}

\subsection*{Your turn}
\begin{itemize}
\raggedright
  \item \emph{Say it:} \emph{die Ausdauer}
  \item \emph{Say it:} \emph{die Ausdauer}, once more
  \item \emph{Say it:} a sentence of your own with \emph{die Ausdauer}
  \item \emph{Recall:} say the German for harvest, then the German for to train, then say \emph{die Ausdauer} again
\end{itemize}

\begin{tcolorbox}[breakable,colback=teal!4,colframe=teal!35!black,title={Before you move on}]
What does die Ausdauer mean? (\textquotedblleft{}Stamina, endurance\textquotedblright{}.) Say it once more.
\end{tcolorbox}
\section[die Kondition]{die Kondition — fitness, physical condition}

\label{lesson:GE-C505-kondition}



\begin{quote}
Before the new one: say the German for to train, then the German for stamina.
\end{quote}

\subsection*{You'll want to know: die Kondition}
\textbf{die Kondition} — \textquotedblleft{}fitness, physical condition\textquotedblright{}.

\begin{quote}
Nach dem Winter ist meine Kondition schlecht.
\end{quote}

\begin{quote}
\emph{After the winter my fitness is poor.}
\end{quote}

No plural in this sense. \emph{Die Konditionen} — terms of a deal.

\begin{cousinweb}
\emph{Kondition} comes from Medieval Latin \emph{conditio}, from Latin \emph{condiciō} \textquotedblleft{}agreement, condition\textquotedblright{} (EtymDB).
\end{cousinweb}

\begin{grammarlens}[title={What you've built}]
One more word for telling, explaining and describing at B1.
\end{grammarlens}

\subsection*{Your turn}
\begin{itemize}
\raggedright
  \item \emph{Say it:} \emph{die Kondition}
  \item \emph{Say it:} \emph{die Kondition}, once more
  \item \emph{Say it:} a sentence of your own with \emph{die Kondition}
  \item \emph{Recall:} say the German for to train, then the German for stamina, then say \emph{die Kondition} again
\end{itemize}

\begin{tcolorbox}[breakable,colback=teal!4,colframe=teal!35!black,title={Before you move on}]
What does die Kondition mean? (\textquotedblleft{}Fitness, physical condition\textquotedblright{}.) Say it once more.
\end{tcolorbox}
\section[dehnen]{dehnen — to stretch}

\label{lesson:GE-C505-dehnen}



\begin{quote}
Before the new one: say the German for stamina, then the German for fitness.
\end{quote}

\subsection*{You'll want to know: dehnen}
\textbf{dehnen} — \textquotedblleft{}to stretch\textquotedblright{}.

\begin{quote}
Vor dem Laufen dehnt sie ihre Beine.
\end{quote}

\begin{quote}
\emph{Before running she stretches her legs.}
\end{quote}

\emph{Sich dehnen} — to stretch yourself.

\begin{cousinweb}
\emph{dehnen} is inherited from Old High German \emph{denen} \textquotedblleft{}to stretch\textquotedblright{} (EtymDB).
\end{cousinweb}

\begin{grammarlens}[title={What you've built}]
One more word for telling, explaining and describing at B1.
\end{grammarlens}

\subsection*{Your turn}
\begin{itemize}
\raggedright
  \item \emph{Say it:} \emph{dehnen}
  \item \emph{Say it:} \emph{dehnen}, once more
  \item \emph{Say it:} a sentence of your own with \emph{dehnen}
  \item \emph{Recall:} say the German for stamina, then the German for fitness, then say \emph{dehnen} again
\end{itemize}

\begin{tcolorbox}[breakable,colback=teal!4,colframe=teal!35!black,title={Before you move on}]
What does dehnen mean? (\textquotedblleft{}To stretch\textquotedblright{}.) Say it once more.
\end{tcolorbox}
\section[der Muskelkater]{der Muskelkater — sore muscles (after exercise)}

\label{lesson:GE-C505-muskelkater}



\begin{quote}
Before the new one: say the German for fitness, then the German for to stretch.
\end{quote}

\subsection*{You'll want to know: der Muskelkater}
\textbf{der Muskelkater} — \textquotedblleft{}sore muscles (after exercise)\textquotedblright{}.

\begin{quote}
Nach dem ersten Training hatte ich drei Tage Muskelkater.
\end{quote}

\begin{quote}
\emph{After my first training session I had sore muscles for three days.}
\end{quote}

No plural. Literally a muscle tomcat; \emph{der Kater} is also a hangover.

\begin{cousinweb}
\emph{Muskelkater} is built on \emph{Muskel}, from Latin \emph{mūsculus} \textquotedblleft{}muscle\textquotedblright{} (EtymDB), and \emph{Kater} \textquotedblleft{}tomcat\textquotedblright{}, inherited from Old High German \emph{kataro}; the joke behind the second part is not settled.
\end{cousinweb}

\begin{grammarlens}[title={What you've built}]
One more word for telling, explaining and describing at B1.
\end{grammarlens}

\subsection*{Your turn}
\begin{itemize}
\raggedright
  \item \emph{Say it:} \emph{der Muskelkater}
  \item \emph{Say it:} \emph{der Muskelkater}, once more
  \item \emph{Say it:} a sentence of your own with \emph{der Muskelkater}
  \item \emph{Recall:} say the German for fitness, then the German for to stretch, then say \emph{der Muskelkater} again
\end{itemize}

\begin{tcolorbox}[breakable,colback=teal!4,colframe=teal!35!black,title={Before you move on}]
What does der Muskelkater mean? (\textquotedblleft{}Sore muscles (after exercise)\textquotedblright{}.) Say it once more.
\end{tcolorbox}
\section[(dialogue)]{Words from chapters 498–505 — a review}

\label{lesson:GE-R505-review-498-505}



\begin{quote}
Say the German for to get to know and on friendly terms.
\end{quote}

\subsection*{Your turn}
Say these aloud:

\begin{itemize}
\raggedright
  \item to get to know, on friendly terms, acquaintance, neighbourhood, hospitality
  \item needle, thread, to knit, to crochet, to embroider
  \item collection, coin, to do crafts, riddle, crossword
  \item board game, chess, dice, to cheat, round
  \item to fish, bait, to bite, lens, to snap
  \item to garden, weeds, to weed, watering can, compost
  \item to dig over, spade, to sow, seed, harvest
  \item to train, stamina, fitness, to stretch, sore muscles
\end{itemize}

\begin{tcolorbox}[breakable,colback=teal!4,colframe=teal!35!black,title={Before you move on}]
To get to know? (\textbf{kennenlernen}.) On friendly terms? (\textbf{befreundet}.) Acquaintance? (\textbf{die Bekanntschaft}.) Neighbourhood? (\textbf{die Nachbarschaft}.) Hospitality? (\textbf{die Gastfreundschaft}.) Needle? (\textbf{die Nadel}.) Thread? (\textbf{der Faden}.) To knit? (\textbf{stricken}.) To crochet? (\textbf{häkeln}.) To embroider? (\textbf{sticken}.) Collection? (\textbf{die Sammlung}.) Coin? (\textbf{die Münze}.) To do crafts? (\textbf{basteln}.) Riddle? (\textbf{das Rätsel}.) Crossword? (\textbf{das Kreuzworträtsel}.) Board game? (\textbf{das Brettspiel}.) Chess? (\textbf{das Schach}.) Dice? (\textbf{der Würfel}.) To cheat? (\textbf{mogeln}.) Round? (\textbf{die Runde}.) To fish? (\textbf{angeln}.) Bait? (\textbf{der Köder}.) To bite? (\textbf{anbeißen}.) Lens? (\textbf{das Objektiv}.) To snap? (\textbf{knipsen}.) To garden? (\textbf{gärtnern}.) Weeds? (\textbf{das Unkraut}.) To weed? (\textbf{jäten}.) Watering can? (\textbf{die Gießkanne}.) Compost? (\textbf{der Kompost}.) To dig over? (\textbf{umgraben}.) Spade? (\textbf{der Spaten}.) To sow? (\textbf{säen}.) Seed? (\textbf{der Samen}.) Harvest? (\textbf{die Ernte}.) To train? (\textbf{trainieren}.) Stamina? (\textbf{die Ausdauer}.) Fitness? (\textbf{die Kondition}.) To stretch? (\textbf{dehnen}.) Sore muscles? (\textbf{der Muskelkater}.)
\end{tcolorbox}
